Long-running LLM conversations inevitably exceed the context window, forcing
systems to compact old messages through summarization. But how do we know
whether our compaction benchmark accurately measures information loss? Before
comparing strategies, we need to understand how reliably an LLM can recall
facts from its own context --- and what factors affect that recall.

We present a three-phase study. First, we \textbf{calibrate a recall benchmark}
using 234 naturally-embedded facts from the LongMemEval dataset across a
190K-token context. We discover four measurement pitfalls that affect all
compaction benchmarks: (1) a \textbf{questions-per-prompt effect} where Q=10 yields
11pp higher recall than Q=1 in static contexts --- but the effect \textbf{reverses}
under severe compaction, with Q=1 outperforming Q=5 by 9pp,
(2) a \textbf{category hierarchy} where temporal reasoning (30\% max) and preference
recall (40\% max) are near-impossible regardless of strategy, (3) a \textbf{density
saturation} where recall plateaus beyond $\approx$0.4 facts/kTok despite increasing
evidence, and (4) a \textbf{grep-LLM gap} where keyword search finds 86\% of facts
but the LLM only recalls 25--79\% of them --- information is \textit{present} in the
context but \textit{ignored} by attention.

Second, we \textbf{isolate compaction loss} by compacting 5--98\% of a context,
re-padding to the original size, and measuring recall degradation. Even 5\%
compaction costs 7 percentage points of recall. At 50\%, the compacted zone is
near-dead (0--7\% recall) despite keywords surviving at 82--93\% via grep.
Critically, compaction damages even the \textit{untouched} portion of the context:
remaining-zone recall drops from 68\% to 39\% as compaction increases --- an
\textbf{attention dilution} effect caused by injecting noise (re-padding) into the
context. Cross-model validation with Claude Sonnet 4.6 (92.5\% baseline recall,
no Lost-in-the-Middle effect) confirms that severe compaction destroys
information regardless of model capability: even a model with flat spatial
recall drops to 21\% at 98\% compaction.

Third, we \textbf{compare four multi-pass compaction strategies} (Brutal,
Incremental, Frozen, FrozenRanked) on a single 5M-token conversation
evaluated at five mid-feed checkpoints (500K $\to$ 5M) with constant fact
density and 4--6 replicates per cell. The strategy hierarchy is consistent
in the means: \textbf{FrozenRanked $>$ Frozen $>$ Incremental $>$ Brutal}. All
strategies degrade severely with scale (Frozen drops from 14.9\% at 500K
to 3.0\% at 5M). With replicates we also document a \textbf{substantial run-to-run
variance}: the compaction phase itself is non-deterministic at temperature
zero and recall measurements on identical conversations span up to a factor
of 14$\times$ (e.g.\ S4 at 1M: 2.6\%--35.9\% across replicates). Single-shot
benchmarks of compaction strategies are therefore unreliable; replicates
are mandatory.

The bottleneck is attention capacity, not compression quality: keywords
survive summarization but the LLM cannot retrieve them, and adding more
preserved summaries dilutes attention rather than helping.

\textit{Scope.} All experiments use Anthropic Claude models (Haiku 4.5 and
Sonnet 4.6). The methodological findings (Q-effect, judge-prompt
sensitivity, run-to-run variance) are likely model-agnostic, but the
absolute recall numbers and the strategy ordering should be re-validated
on non-Claude models before generalising.
